% Shear stress vs. time for the torsion problem with RaM, one CSV per RaM segment (_0 to _114).
% Columns are referenced by name because later segments add a dt column.
% Segments _1 onward start with an all-zero row at the restart time, so that row is dropped.
\pgfplotsset{
    torskipfirst/.style={x filter/.expression={\coordindex==0 ? nan : x}},
}
\begin{figure}[htbp]
    \centering
    \begin{tikzpicture}
        \begin{axis}[
                width=0.6\textwidth,
                height=0.4\textwidth,
                xmin=0,
                xlabel={Time (\unit{\micro\second})},
                ylabel={Shear stress (\unit{\mega\pascal})},
                grid=major,
            ]
            \addplot[blue, thick] table[col sep=comma, x=time, y=shear_stress] {results/data/torsion/ss/tor_d_small_1xs_new_0.csv};
            \foreach \i in {1,...,114} {
                    \addplot[blue, thick, torskipfirst] table[col sep=comma, x=time, y=shear_stress] {results/data/torsion/ss/tor_d_small_1xs_new_\i.csv};
                }
        \end{axis}
    \end{tikzpicture}
    \caption{Shear stress vs.\ time for the torsion problem with RaM.}\label{fig:torsion_shear_stress}
\end{figure}
